\section{Specification}


\subsection{Components}

The replication system consists of three components:

\begin{table}[H]
\centering
\begin{tabular}{@{}lll@{}}
\toprule
\textbf{Component} & \textbf{Image} & \textbf{Purpose} \\
\midrule
\texttt{hanzoai/replicate} & \texttt{ghcr.io/hanzoai/replicate} & WAL replication + age encryption \\
\texttt{hanzoai/age}       & \texttt{ghcr.io/hanzoai/age}       & X25519 key generation \\
\texttt{hanzoai/s3}        & \texttt{ghcr.io/hanzoai/s3}        & S3-compatible object storage \\
\bottomrule
\end{tabular}
\caption{System components.}
\label{tab:components}
\end{table}

\texttt{hanzoai/replicate} is a fork of Litestream incorporating Beeper's
age encryption patches, rebased onto mainline Litestream, and maintained as
part of the Hanzo ecosystem. The binary is called \texttt{replicate}. It
accepts the same configuration format as Litestream with one addition: an
\texttt{age} block specifying identities and recipients.

\texttt{hanzoai/age} is a container image packaging the \texttt{age-keygen}
binary for key generation in init containers and CI pipelines.

\texttt{hanzoai/s3} (HIP-0032) provides the S3-compatible storage backend.

\subsection{Replication Protocol}

The replication lifecycle for a single SQLite database:

\begin{enumerate}[nosep]
    \item \textbf{Startup.} \texttt{replicate} reads the configuration file,
    connects to the S3 endpoint, and checks for an existing replica.

    \item \textbf{Restore (if needed).} If the local database does not exist
    and a replica exists in S3, \texttt{replicate} downloads the most recent
    snapshot, decrypts it with the age identity, replays all subsequent WAL
    segments (each decrypted individually), and writes the restored database
    to the configured path.

    \item \textbf{Monitor.} \texttt{replicate} monitors the SQLite WAL file
    using \texttt{inotify} (Linux) or \texttt{kqueue} (macOS). When new
    frames are written, it reads them from the WAL.

    \item \textbf{Encrypt.} Each WAL segment is encrypted with the
    configured age recipients before upload. Encryption happens in memory.
    Plaintext WAL data never touches S3.

    \item \textbf{Upload.} The encrypted WAL segment is uploaded to S3 at the
    configured path. The upload is atomic (multipart with commit).

    \item \textbf{Snapshot.} Periodically (default: every 24 hours),
    \texttt{replicate} takes a full snapshot of the database, encrypts it,
    and uploads it to S3. Old snapshots are retained according to the
    configured retention policy.

    \item \textbf{Shutdown.} On SIGTERM, \texttt{replicate} flushes any
    pending WAL frames, uploads them, and exits cleanly. Kubernetes sends
    SIGTERM before SIGKILL, so graceful shutdown is the normal path.
\end{enumerate}

\subsection{Encryption Layers}

Data is protected by three independent encryption layers:

\begin{table}[H]
\centering
\begin{tabular}{@{}llll@{}}
\toprule
\textbf{Layer} & \textbf{Mechanism} & \textbf{Scope} & \textbf{Protects Against} \\
\midrule
1. Disk   & sqlcipher (AES-256-CBC)      & Database at rest     & Local filesystem access \\
2. S3     & age (X25519 + ChaCha20)      & Replicas in storage  & S3 operator, bucket leak \\
3. Wire   & TLS 1.3                       & In transit           & Network interception \\
\bottomrule
\end{tabular}
\caption{Encryption layers. Each layer is independent---compromise of one
does not compromise the others.}
\label{tab:encryption}
\end{table}

Layer 1 (sqlcipher) is managed by Hanzo Base. Layer 2 (age) is managed by
\texttt{replicate}. Layer 3 (TLS) is managed by the S3 endpoint. This HIP
specifies Layer 2.

\subsection{Recovery Objectives}

\begin{table}[H]
\centering
\begin{tabular}{@{}lll@{}}
\toprule
\textbf{Metric} & \textbf{Target} & \textbf{Mechanism} \\
\midrule
RPO (Recovery Point Objective) & $\leq$ 1 second & Continuous WAL streaming \\
RTO (Recovery Time Objective)  & $\leq$ 30 seconds & Snapshot restore + WAL replay \\
Retention                      & 72 hours WAL, 30 days snapshots & S3 lifecycle policy \\
\bottomrule
\end{tabular}
\caption{Recovery objectives.}
\label{tab:recovery}
\end{table}

The 1-second RPO is achieved by the default \texttt{sync-interval} of 1
second. WAL frames generated within the last sync interval may be lost on
unclean shutdown (SIGKILL without preceding SIGTERM). In practice,
Kubernetes provides a 30-second grace period between SIGTERM and SIGKILL, so
data loss on clean shutdown is zero.

The 30-second RTO accounts for S3 download latency for the most recent
snapshot (typically 1--5 seconds for databases under 1~GiB) plus WAL replay
time (linear in the number of frames since the last snapshot).

% ============================================================================
% 4. CONFIGURATION
% ============================================================================
